\begin{center}
\textbf{\large{Abstract}	}
\end{center}

\addcontentsline{toc}{chapter}{Abstract}
\begin{small}

% Summary of the graduation thesis: The problem statement, the proposed approach, and the achieved results.

This thesis investigates the application of machine learning to the
detection of Portable Executable (PE) malware on the Windows operating
system, with two main contributions: an empirical evaluation of
capability extraction as a feature engineering technique in machine
learning pipelines, and the design and implementation of a functional
endpoint protection system on Windows.

On the research side, three LightGBM-based machine learning pipelines
are proposed and compared using the EMBER2024 dataset. The pipelines
differ in how their input feature vectors are constructed: using only
EMBER2024 features (baseline), augmenting with Capa-derived capability
features encoded as one-hot vectors, and augmenting with Capa features
reduced via TruncatedSVD. Experiments conducted on 450,000 training
samples and 50,000 test samples show that the pipeline combining
EMBER2024 with Capa TruncatedSVD-256 achieves the best performance,
with a Partial AUC (FPR $\leq$ 0.5\%) of 0.8850 and a Full AUC of
0.9926, improving over the baseline by 0.34\% and 0.04\% respectively.
These results demonstrate that the behavioral semantic information
provided by Capa contributes meaningful discriminative value beyond
conventional PE structural features.

On the engineering side, the thesis designs and implements a five-
component system architecture: a high-performance C++ feature
extraction module (12 to 20 times faster than the original Python
implementation of the EMBER2024 project), a KMDF-based kernel driver for process execution
monitoring, a Windows background service integrating the machine
learning model, a real-time notification user interface, and an online
deep analysis service incorporating Capa and SHAP. The system is
capable of detecting and blocking malware at the moment of execution
on the endpoint, while also providing in-depth analysis and model
decision explanations through the online deep analysis service.

\vspace*{1cm}
\textbf{Keywords}: malware detection, machine learning,
EMBER2024, capability extraction, Windows driver.
\\
\textbf{GitHub repository}: \url{https://github.com/laam-egg/KLTN_22028235}

\end{small}
